\loai{Khối lượng phân rã}
\begin{vidu}%[3P7-7.4B][Loai3] 
\nguon{QG - 2019} Chất phóng xạ X có chu kì bán rã là T. Ban đầu có một mẫu X nguyên chất với khối lượng 4 g. Sau khoảng thời gian 2T, khối lượng chất X trong mẫu đã bị phân rã là
\choice
{1 g}
{\True 3 g}
{2 g}
{0,25 g}
\loigiai{
Khối lượng phân rã: $\Delta m=4-\dfrac{4}{2^2}=3\ g$}
\end{vidu}
\loai{Số nguyên tử phân rã}
\begin{vidu}%[3P7-7.4B][Loai2]
Pôlôni \Hn{84}{210}{Po} là chất phóng xạ có chu kì bán rã 138 ngày. Lấy $N_A=6,023.10^{23}\ mol^{-1}$. Lúc đầu có 10 g polonium.  Lấy khối lượng nguyên tử tính theo đơn vị u bằng số khối của hạt nhân của nguyên tử đó. Sau thời gian 69 ngày, số nguyên tử polonium bị phân rã là
\choice
{\True $8,4.10^{21}$}
{$6,5.10^{22}$}
{$2,9.10^{20}$}
{$5,7.10^{23}$}	
\loigiai{
Sau 69 ngày, số nguyên tử polonium  bị phân rã là:
\begin{align*}
N=N_0\pbrace{1-2^{-\tfrac{t}{T}}} = \dfrac{m.N_A}{A}\pbrace{1-2^{-\tfrac{t}{T}}} = \dfrac{10.6,{023.10}^{23}}{210}.\pbrace{1-2^{-69/138}}=8,4.10^{21}\text{hạt}.
\end{align*}
}
\end{vidu}
\loai{Số hạt phóng xạ}
\begin{vidu}%[3P7-7.4K][Loai5]
\Hn{27}{60}{Co} là chất phóng xạ $\beta^-$ với chu kì bán rã 5,27 năm.  Lấy khối lượng nguyên tử tính theo đơn vị u bằng số khối của hạt nhân của nguyên tử đó. Tính số hạt $\beta^-$ phát ra từ 0,6 g \Hn{27}{60}{Co} sau 15,81 năm.
Cho $N_A=6,02.10^{23}$ hạt/mol.	
\choice
{\True $5,2675.10^{21}$ hạt}
{$5,2675.10^{22}$ hạt}
{$7,525.10^{20}$ hạt}
{$7,525.10^{21}$ hạt}
\loigiai{
\begin{itemize}
\item Số hạt Co có trong 0,6 g là: $N_0=N_A.\dfrac{m_{Co}}{M_{Co}}=6,02.10^{23}.\dfrac{0,6}{60}=6,02.10^{21}$ hạt.
\item Số hạt $\beta^-$ phát ra sau 15,81 năm bằng số hạt phân rã: 
\begin{align*}
\Delta N=N_0\pbrace{1 - 2^{-\tfrac{t}{T}}}=6,02.10^{21}\pbrace{1 - 2^{ - \tfrac{15,81}{5,27}}}=5,2675.10^{21}\ \text{hạt}.
\end{align*}
\end{itemize}
}
\end{vidu}
\loai{Số hạt tạo thành}
\begin{vidu}%[3P7-7.4K][Loai4]
Đồng vị phóng xạ $\Hn{88}{226}{Ra}$ phân rã $\alpha$ và biến đổi thành hạt nhân X. Lúc đầu $\Hn{88}{226}{Ra}$ nguyên chất có khối lượng 0,064 g. Hạt nhân Ra có chu kì bán rã là 1517 năm. Lấy khối lượng nguyên tử tính theo đơn vị u bằng số khối của hạt nhân của nguyên tử đó. Số hạt nhân X tạo thành trong năm thứ 786 là
\choice
{\True $5,44.10^{16}$ hạt}
{$4,57.10^{15}$ hạt}
{$4.10^{16}$ hạt}
{$2,28.10^{16}$ hạt}
\loigiai{
\begin{itemize}
\item Số hạt nhân Ra lúc đầu là: ${N_0}=\dfrac{m}{M}. {N_A}=\dfrac{0,064}{226}. 6,02.10^{23}=1,705.10^{20}$ hạt.
\item Số hạt nhân Ra còn lại sau ${t_1}=785$ năm là: 
\begin{align*}
N_{t_1}=N_0. 2^{ - \tfrac{t_1}{T}}=1,705.10^{20}.2^{ - \tfrac{785}{1517}}=1,19095.10^{20}\ \text{hạt}.
\end{align*}
\item Số hạt nhân Ra còn lại sau ${t_2}=786$ năm là: 
\begin{align*}
N_{t_2}=N_0. 2^{ - \tfrac{t_2}{T}}=1,705.10^{20.2^{ - \tfrac{786}{1517}}}=1,19041.10^{20}\ \text{hạt}.
\end{align*}
\item Số hạt nhân X tạo thành trong năm thứ 786 bằng số hạt nhân Ra bị phân rã trong năm thứ 786 
\begin{align*}
{N_X}=N_{t_1}-N_{t_2}=5,44.10^{16}\ \ct{hạt}.
\end{align*}
\end{itemize}
}
\end{vidu}
\loai{Tỉ số số hạt mẹ và con}
\begin{vidu}%[3P7-7.4K][Loai6]
\nguon{MH - 2018} Hạt nhân X phóng xạ biến đổi thành hạt nhân bền Y. Ban đầu $\left(t=0\right)$ có một mẫu chất X nguyên chất. Tại thời điểm $t_1$ và $t_2$ tỉ số giữa số hạt nhân Y và số hạt nhân X ở trong mẫu tương ứng là 2 và 3. Tại thời điểm $t_3=2t_1+3t_2$, tỉ số đó là
\choice
{17}
{\True 575}
{107}
{72}
\loigiai{
\begin{itemize}
\item Ta có $\begin{cases}
&N_Y=N_0\left(1-2^{-\tfrac{t}{T}}\right)\\
&N_X=N_02^{-\tfrac{t}{T}}\\
\end{cases}$
\item Thời điểm $t_1:\ \dfrac{N_Y}{N_X}=\dfrac{1-2^{-\tfrac{t_1}{T}}}{2^{-\tfrac{t_1}{T}}}=2\Rightarrow 2^{-\tfrac{t_1}{T}}=\dfrac{1}{3}\quad (1)$.
\item Thời điểm $t_2:\ \dfrac{N_Y}{N_X}=\dfrac{1-2^{-\tfrac{t_2}{T}}}{2^{-\tfrac{t_2}{T}}}=2\Rightarrow 2^{-\tfrac{t_2}{T}}=\dfrac{1}{4}\quad (2)$.
\item Thời điểm $t_3=2t_1+3t_2:\ \dfrac{N_Y}{N_X}=\dfrac{1-2^{-\tfrac{2t_1+3t_2}{T}}}{2^{-\tfrac{2t_1+3t_2}{T}}}=\dfrac{1-2^{-\tfrac{2t_1}{T}}.2^{-\tfrac{3t_2}{T}}}{2^{-\tfrac{2t_1}{T}}.2^{-\tfrac{3t_2}{T}}}=\dfrac{1-{\left(2^{-\tfrac{t_1}{T}}\right)}^2.{\left(2^{-\tfrac{t_2}{T}}\right)}^3}{{\left(2^{-\tfrac{t_1}{T}}\right)}^2.{\left(2^{-\tfrac{t_2}{T}}\right)}^3}\quad (3)$.
\item Thay $(1)$ và $(2)$ vào $(3)$ ta được: $\dfrac{N_Y}{N_X}=575$.
\end{itemize}
}
\end{vidu}
\loai{Khối lượng chất tạo thành}
\begin{vidu}%[3P7-7.4G][Loai7]
\nguon{QG - 2018} Chất phóng xạ poloni $\Hn{84}{210}{Po}$ phát ra tia alpha và biến đổi thành chì $\Hn{82}{206}{Pb}$. Gọi chu kì bán rã của poloni là T. Ban đầu ($t = 0$) có một mẫu $\Hn{84}{210}{Po}$ nguyên chất. Trong khoảng thời gian từ $t = 0$ đến $t = 2T$ có 126 mg $\Hn{84}{210}{Po}$ trong mẫu bị phân rã. Lấy khối lượng nguyên tử tính theo đơn vị u bằng số khối của hạt nhân của nguyên tử đó. Trong khoảng thời gian từ $t = 2T$ đến $t = 3T$, lượng $\Hn{82}{206}{Pb}$ được tạo thành trong mẫu có khối lượng là
\choice
{61,8 mg}
{41,2 mg}
{\True 20,6 mg}
{10,5 mg}
\loigiai{
\begin{itemize}
\item Khối lượng chất phóng xạ bị phân rã từ $t = 0$ đến $t = 2T$:
\begin{align*}
m_0 - \dfrac{m_0}{2^2}=126\ mg\Rightarrow m_0=168\ mg.
\end{align*}
\item Khi đó, thời điểm $t=2T$, khối lượng Po còn lại là: $m_0'=\dfrac{m_0}{4}=42\ mg$.
\item  Số hạt Po bị phân rã trong thời gian từ 2T đến 3T là 
\begin{align*}
N_{Pb}=\Delta N_{Po}=N_{0Po}.2^{-\tfrac{t}{T}}=\dfrac{m_0'}{A_{Po}}.N_A.2^{ - \tfrac{3T - 2T}{T}}\Rightarrow m_{Pb}=\dfrac{m_0'}{A_{Po}}.A_{Pb}=\dfrac{42}{210}.2^{-1}.206=20,6\ mg.
\end{align*}
\end{itemize}
}
\end{vidu}

\loai{Tỉ số khối lượng mẹ và con}
\begin{vidu} %[3P7-7.4G][Loai8]
Hạt nhân $\Hn{{Z_1}}{{A_1}}{X}$ phân rã và trở thành hạt nhân $\Hn{{Z_2}}{{A_2}}{Y}$ bền. Coi khối lượng hai hạt nhân đó bằng số khối của chúng tính theo đơn vị u. Lúc đầu mẫu $\Hn{{Z_1}}{{A_1}}{X}$ là nguyên chất. Biết chu kì phóng xạ của $\Hn{{Z_1}}{{A_1}}{X}$ là T (ngày). Ở thời điểm $T + 14$ (ngày) tỉ số khối lượng của $\Hn{{Z_1}}{{A_1}}{X}$ và $\Hn{{Z_2}}{{A_2}}{Y}$ là $\dfrac{A_1}{7{A_2}}$; đến thời điểm $T + 28$ (ngày) tỉ số khối lượng trên là
\choice
{$\dfrac{A_1}{14{A_2}}$}
{$\dfrac{7{A_1}}{8{A_2}}$}
{\True $\dfrac{A_1}{31{A_2}}$}
{$\dfrac{A_1}{32{A_2}}$}
\loigiai{
\begin{itemize}
\item Ở thời điểm $T + 14$ ngày: 
\begin{align*}
\dfrac{m_Y}{m_X}=\dfrac{N_Y}{N_X}.\dfrac{A_2}{A_1}=\dfrac{7A_2}{A_1}\Rightarrow\left({2^{\tfrac{T + 14}{T}}} - 1\right)=7\Rightarrow{{2. 2}^{\tfrac{14}{T}}} - 1=7\Rightarrow\dfrac{14}{T}=2\Rightarrow T=7\ \text{ngày}.
\end{align*}
\item Đến thời điểm $T + 28$ ngày:
\begin{align*}
\dfrac{m'_X}{m'_Y}=\dfrac{N'_X}{N'_Y}. \dfrac{A_1}{A_2}=\dfrac{1}{
{2^{\tfrac{T+28}{T}}}-1}\dfrac{A_1}{A_2}=\dfrac{1}{{2^{\tfrac{7+28}{7}}}-1}\dfrac{A_1}{A_2}=\dfrac{A_1}{31{A_2}}.
\end{align*}
\end{itemize}
}	
\end{vidu}
\loai{Tính khối lượng ban đầu}
\begin{vidu}%[3P7-7.4K][Loai9]
Đồng vị \Hn{92}{238}{U} sau một chuỗi các phân rã thì biến thành chì \Hn{82}{206}{Pb} bền, với chu kì bán rã $T=4,47$ tỉ năm. Ban đầu có một mẫu chất $^{238}U$ nguyên chất. Sau 2 tỉ năm thì trong mẫu chất có lẫn chì $^{206}Pb$ với khối lượng $m_{Pb}=0,2\ g$. Giả sử toàn bộ lượng chì đó đều là sản phẩm phân rã từ $^{238}U$. Khối lượng $^{238}U$ ban đầu là
\choice
{0,428 g}
{4,28 g}
{\True 0,866 g}
{8,66 g}
\loigiai{
Khối lượng chì được tạo thành sau 2 tỉ năm: $m_{Pb}=\dfrac{206}{238}m_0\left(1-2^{-\tfrac{2}{4,47}}\right)\Rightarrow m_0=0,866\ g$.
}
\end{vidu}
\loai{Tính thời gian}
\begin{vidu}%[3P7-7.4G][Loai10]
\nguon{QG - 2017} Chất phóng xạ polonium  \Hn{84}{210}{Po} phát ra tia $\alpha$ biến đổi thành hạt nhân chì. Chu kì bán rã của polonium  là 138 ngày. Ban đầu có một mẫu polonium  nguyên chất, sau khoảng thời gian t, tỉ số giữa khối lượng chì sinh ra và khối lượng polonium  còn lại trong mẫu là 0,6. Coi khối lượng nguyên tử bằng số khối của hạt nhân của nguyên tử đó tính theo đơn vị u. Giá trị của t là	
\choice
{\True 95 ngày}
{83 ngày}
{33 ngày}
{105 ngày}	
\loigiai{
\begin{itemize}
\item Khối lượng Po còn lại trong mẫu: $m_{Po}=m_{0}.2^{-\tfrac{t}{T}}=\dfrac{N_0.N_{Pb}}{N_A}$.
\item Số hạt nhân Po bị phân rã chính là số hạt nhân Pb được sinh ra: 
\begin{align*}
\Delta N=N_{Pb}=N_0-N=N_0\left( 1-2^{-\tfrac{t}{T}}\right).
\end{align*}
\item Khối lượng Pb sinh ra: $m_{Pb}=\dfrac{N_{Pb}.A_{Pb}}{N_A}=\dfrac{N_0\left( 1-2^{-\tfrac{t}{T}}\right).A_{Pb}}{N_A}$.
\item Tỉ số giữa khối lượng Pb sinh ra và khối lượng Po còn lại trong mẫu: 
\begin{align*}
\dfrac{m_{Pb}}{m_{Po}}=\dfrac{206}{210}\dfrac{\pbrace{1-2^{-\tfrac{t}{138}}}}{{2^{-\tfrac{t}{138}}}}=0,6\Rightarrow t=95\ \text{ngày}.
\end{align*}
\end{itemize}	
}
\end{vidu}
\loai{Tính chu kì bán rã}
\begin{vidu}%[3P7-7.4K][Loai12]
Chất phóng xạ $\Hn{84}{210}{Po}$ phóng xạ $\alpha$ rồi trở thành chì (Pb). Dùng một mẫu Po ban đầu có 1 g, sau 365 ngày đêm mẫu phóng xạ trên tạo ra lượng khí helium có thể tích là: $V = 89,5\ cm^3$ ở điều kiện tiêu chuẩn. Lấy khối lượng nguyên tử tính theo đơn vị u bằng số khối của hạt nhân của nguyên tử đó. Cho $N_A=6,023.10^{23}\ mol^{-1}$. Chu kì bán rã của Po là
\choice
{\True 138,5 ngày đêm}
{135,6 ngày đêm}
{148 ngày đêm}
{138 ngày đêm}
\loigiai{
\begin{itemize}
\item Số hạt khí helium tạo thành là: $N_{He}=n_{He}. {N_A}=\dfrac{0,0895}{22,4}. 6,02.10^{23}=2,405.10^{21}$ hạt.
\item Ta có: $N_{He}=\Delta N\Rightarrow N_{He}={N_0}\left( 1-{2^{-\tfrac{t}{T}}} \right)\Rightarrow 2,405.10^{21}=\dfrac{1}{210}. 6,02.10^{23}. \left( 1-{2^{-\tfrac{365}{T}}} \right)\Rightarrow T=138,5$ ngày đêm.
\end{itemize}
}
\end{vidu}

\loai{Hai chất phóng xạ}
\begin{vidu}%[3P7-7.4K][Loai13]
Chu kì bán rã của hai chất phóng xạ A và B lần lượt là $T_A$ và $T_B=2T_A$. Ban đầu hai khối chất A và B có số hạt nhân như nhau. Sau thời gian $t=4T_A$ thì tỉ số giữa số hạt nhân A và B đã phóng xạ là 
\choice
{$\dfrac{1}{4}$}
{4}
{$\dfrac{4}{5}$}
{\True $\dfrac{5}{4}$}
\loigiai{
Số hạt nhân đã phóng xạ:
\begin{align*}
\Delta N=N_0\left(1-2^{-\tfrac{t}{T}}\right)\Rightarrow  
\begin{cases}
\Delta N_1=N_0\left({1-2^{-\tfrac{t}{T_A}}}\right) \\
\Delta N_2=N_0\left({1-2^{-\tfrac{t}{2T_A}}}\right) 
\end{cases}\Rightarrow \dfrac{\Delta N_1}{\Delta N_2}=\dfrac{1-2^{-4}}{1-2^{-2}}=\dfrac{5}{4}.
\end{align*}
}
\end{vidu}
\loai{Bài toán tổng hợp}
\begin{vidu}%[3P7-7.4G][Loai14]
Chất polonium \Hn{84}{210}{Po} là là phóng xạ hạt $^4\alpha $ có chu kì bán rã là 138 ngày. Ban đầu giả sử mẫu quặng $Po$ là nguyên chất và có khối lượng $210\ g$. Chì được sinh ra nằm ở trong mẫu. Lấy khối lượng nguyên tử tính theo đơn vị u bằng số khối của hạt nhân của nguyên tử đó. Sau 276 ngày người ta đem mẫu quặng đó ra cân thì khối lượng còn lại của mẫu quặng là
\choice
{157,5 g}
{52,5 g}
{210 g}
{\True 207 g}
\loigiai{
\begin{itemize}
\item Phương trình phóng xạ: $\Hn{84}{210}{Po}\to^4\alpha +^{206}X$. 
\item Ta thấy khối lượng của mẫu quặng sẽ bằng tổng khối lượng còn lại của Po và khối lượng của X (hạt $\alpha$ phân tán ra ngoài môi trường dưới dạng tia phóng xạ).
\item Khối lượng Po còn lại trong mẫu: $m_{Po}=m_0.2^{-\tfrac{t}{T}}=210.2^{-\tfrac{276}{138}}=52,5\ g$.
\item Khối lượng X (chì) sinh ra (còn lưu lại hoàn toàn trong mẫu): 
\begin{align*}
m_{Pb}=\Delta m_{Po}.\dfrac{206}{210}=\left(m_0-m \right).\dfrac{206}{210} = \left(210-52,5 \right).\dfrac{206}{210}=154,5\ g.
\end{align*}
\item Vậy khối lượng của mẫu quặng:  $m=m_{Po}+m_{Pb}=52,5+154,5=207\ g$.
\end{itemize}
}
\end{vidu}
